%! Boxplots, Outliers, and Resistance — Unit 2, Lesson 2.3 — Lesson Plan
% PILOT SHAPE (2026-09-06, modeled on AP Statistics 1.4 minus AP Practice), on the UNIT 2 PERIOD
% (2026-09-29): 5 / 45 / 5 / 5 — warm-up / guided notes & practice (I do ~25 row by row, we do
% the Guided Practice ~20) / SPOKEN debrief (cold check + one board item) / close and ASSIGN the
% homework (not started in class). There is no independent practice set in the notes — the
% homework is the individual practice. Teacher-facing: the formal vocabulary and the lettered
% standard codes are used freely. Breakable boxes are `skillbox` + \boxguard; the two tabularx
% boxes are `fixedskillbox` (never splits). Mirrors unit01/lesson02.
\documentclass[10pt]{article}
\usepackage{saar-article}
\usepackage{saar-boxes}
\usepackage{graphicx}   % -article does NOT load graphicx; needed for thumbnails/screenshots
\graphicspath{{images/}}
\newcommand{\TallMath}[1]{$\displaystyle #1\rule[-1.4em]{0pt}{3.2em}$}
\newcommand{\UnitNumberName}{Unit 2: Describing One Variable \quad}
\newcommand{\LessonNumberName}{Lesson 2.3: Boxplots, Outliers, and Resistance}

\begin{document}
\begin{center}
    {\Large\bfseries \CourseName \\
    {\small\bfseries \UnitNumberName \LessonNumberName}}
\end{center}

\begin{tcolorbox}[colback=frost, colframe=cerulean, boxrule=0.9pt, arc=2mm,
    left=2mm, right=2mm, top=1.5mm, bottom=1.5mm]
    \textbf{Primary Objective:} Students will read a boxplot and finish one from a five-number
    summary, use the $1.5 \times \text{IQR}$ rule to find fences and name outliers (and read a
    modified boxplot), read quartiles and percentiles off a cumulative relative frequency graph
    and say them in context, and \emph{justify} --- by computing each summary with and without
    an outlier --- why the median and IQR are the center and spread to report when data have
    outliers or strong skew.
    \\[2pt]
    % CITE THE STANDARD CODES. They appear here and again on the Connections line.
    \textbf{Standards:} PS.DS.1a (boxplots and cumulative frequency graphs --- create and
    interpret) \quad$\cdot$\quad PS.DS.2c (identify possible outliers using an algorithm)
    \quad$\cdot$\quad PS.DS.2d (the influence of outliers) \quad$\cdot$\quad builds on
    Lesson~2.2's five-number summary, IQR, and standard deviation
    \\[2pt]
    \textbf{Lesson model:} \emph{Gradual release --- I do, we do, you do --- all of it inside
    the guided notes.} The notes deliver the instruction \textbf{row by row in one Main Ideas /
    Notes table} (four rows, problems 1--18: you read the definition, mark up the display, and
    work the first problem of each grid --- 1, 5, 9, 13; the class works the rest); the Guided
    Practice row (problems 19--22) is worked together with students holding the pen;
    \textbf{the homework is the individual practice}, assigned at the close and done outside
    class. The debrief is a \emph{spoken} cold check.
    \\[2pt]
    \textbf{Conventions on the page:} quartiles are the medians of the lower and upper halves,
    with the overall median left out of both halves when $n$ is odd (printed once, in notes
    row~1); every standard deviation is the \emph{sample} standard deviation ($\div\,(n-1)$),
    given as calculator output --- students do not compute one today. Single-group boxplots
    only: parallel boxplots are Lesson~2.4.
\end{tcolorbox}

\renewcommand{\arraystretch}{1.4}
\boxguard
\begin{skillbox}[Learning Targets \& Key Understandings]{goldbox}
\begin{tabularx}{\linewidth}{@{}X|X@{}}
\textbf{Learning targets (student language)} & \textbf{Key understandings (the ``why'')} \\[2pt]
\hline
\vspace{-6pt}
\begin{itemize}[leftmargin=1.1em, nosep, after=\vspace{2pt}]
  \item I can read a \textbf{boxplot} and finish one from a five-number summary, knowing each
        of its four sections holds about $25\%$ of the data. \hfill \textit{PS.DS.1a}
  \item I can use the \textbf{$1.5 \times \text{IQR}$ rule} to find the fences, name every
        \textbf{outlier}, and read a \textbf{modified boxplot}. \hfill \textit{PS.DS.2c}
  \item I can read the median, the quartiles, and a \textbf{percentile} off a
        \textbf{cumulative relative frequency graph} and say it in context.
        \hfill \textit{PS.DS.1a}
  \item I can show, with and without an outlier, that the median and IQR are
        \textbf{resistant} and the mean and standard deviation are not --- and choose what to
        report. \hfill \textit{PS.DS.2d}
\end{itemize}
&
\vspace{-6pt}
\begin{itemize}[leftmargin=1.1em, nosep, after=\vspace{2pt}]
  \item \textbf{Target misconception:} that an outlier changes every summary equally --- that
        the median is pulled like the mean. It is not: removing Pine Street's $55$ drops the
        mean $3$ minutes and the median half a minute, and leaves the IQR at $5$. The mean and
        SD use every \emph{value}; the median and IQR use only \emph{positions}.
        \emph{Median and IQR resist; mean and SD do not.}
  \item \textbf{Second misconception:} calling a point an outlier because it ``looks far.''
        The test is the fence, not the eye: a $16$-minute delivery sits inside the lower fence
        of $15.5$ (problem~7); Elm Park's \$40 bike sits inside its fence of \$30 (HW~2).
  \item Each section of a boxplot holds about a quarter of the data; a long whisker means
        spread-out values, not more of them.
  \item A cumulative relative frequency graph is read \emph{across then down}: the value at
        $p\%$ is the $p$th percentile, so $Q_1$, the median, and $Q_3$ are the $25$th, $50$th,
        and $75$th.
  \item The choice of summary is a claim about the data: with outliers or strong skew, the
        mean describes no typical case, so report the median and IQR (PS.DS.2d).
\end{itemize}
\end{tabularx}
\end{skillbox}

\boxguard[24]
\begin{skillbox}[{Vocabulary, Concepts, \& Theorems --- taught directly in the guided notes}]{greenbox}
\small
\renewcommand{\arraystretch}{1.35}
\begin{tabularx}{\linewidth}{@{} >{\bfseries}p{3.4cm} X @{}}
Five-number summary & Minimum, $Q_1$, median, $Q_3$, maximum (Lesson~2.2). \\
Boxplot & A graph of the five-number summary: a box from $Q_1$ to $Q_3$ with a line at the
       median, whiskers out to the extremes. Each of the four sections holds about $25\%$. \\
Outlier ($1.5 \times$ IQR rule) & A value below $Q_1 - 1.5 \times \text{IQR}$ or above
       $Q_3 + 1.5 \times \text{IQR}$ --- the \textbf{lower} and \textbf{upper fences}. \\
Modified boxplot & A boxplot that plots each outlier as its own dot; the whisker stops at the
       most extreme value that is not an outlier. \\
Cumulative relative frequency graph & Plots the percent of the data at or below each value
       (an \emph{ogive}). Read across from $25\%$, $50\%$, $75\%$ for $Q_1$, median, $Q_3$. \\
Percentile & The $p$th percentile is the value with $p\%$ of the data at or below it. \\
Resistant & A summary that one extreme value barely moves. Median and IQR: resistant. Mean and
       standard deviation: not. \\
\end{tabularx}
\end{skillbox}

% fixedskillbox never splits, so the phase table stays intact on one page.
% THE PHASE TABLE IS A CONTRACT: 5 / 45 / 5 / 5 (Unit 2 on; 2026-09-29) — I do ~25 of the 45,
% Guided Practice ~20; the debrief is the spoken cold check; the homework is assigned, not started.
\begin{fixedskillbox}[Lesson at a Glance --- \MeetingLength]{frost}
\renewcommand{\arraystretch}{1.25}
\begin{tabularx}{\linewidth}{@{}l c X X@{}}
\textbf{Phase} & \textbf{Min} & \textbf{Students} & \textbf{Teacher} \\
\hline
Warm-Up & 5 & Order Pine Street's ten Friday deliveries; the median, $Q_1$, $Q_3$, IQR; three
  counts as percents of $40$ --- silent & Debrief item~2 aloud; write
  \emph{$Q_1 = 23$, median $25.5$, $Q_3 = 28$, IQR $5$} on the board and leave it up all
  period \\
Guided Notes \& Practice & 45 & Fill the vocabulary box as each term is named; vote on the
  hook; work rows 1--4 (problems 1--18) with the pen in their hand; then the Guided Practice
  row, \emph{Shelter Stays} (19--22), together & \textbf{I do / we do, row by row}: read the
  definition, mark up the display, work 1, 5, 9, 13; the class works the rest (25) $\to$
  \textbf{we do} Guided Practice 19--22 (20) \\
Debrief & 5 & Pens down --- answer aloud, nothing to fill in & Put the with/without board back
  up and have a student say the sentence; then the cold check \\
Close \& Assign & 5 & Write down the homework, due date, and first item & Announce the
  homework and its form (DeltaMath + packet items 2, 4, 6), 12 points, due after two study halls; preview Lesson~2.4 \\
\end{tabularx}
\end{fixedskillbox}

\begin{fixedskillbox}[Warm-Up --- Activate Prior Knowledge (5 min)]{frost}
\begin{tabularx}{\linewidth}{@{}X X@{}}
\begin{minipage}[t]{\linewidth}\vspace{0pt}
    \textbf{The three items and what each seeds}
    \begin{itemize}[leftmargin=1.1em, itemsep=1pt, topsep=2pt]
      \item \textbf{1.} Order ten delivery times; the median ($25.5$). \textbf{Spirals} to
            Lesson~2.2 --- and it is the list the hook prints in order, and the median problem~14
            compares to $25$.
      \item \textbf{2.} $Q_1 = 23$, $Q_3 = 28$, IQR $= 5$ by the halves method. \textbf{Seeds:}
            the boxplot in row~1 (problem~1 reads these back) and the fences in row~2
            (problem~5 starts from them).
      \item \textbf{3.} $10$, $20$, $30$ of $40$ as $25\%$, $50\%$, $75\%$. \textbf{Seeds:} the
            cumulative relative frequency graph in row~3 --- those are its three quartile
            points.
    \end{itemize}
\end{minipage}
&
\begin{minipage}[t]{\linewidth}\vspace{0pt}
    \textbf{Running it}
    \begin{itemize}[leftmargin=1.1em, itemsep=1pt, topsep=2pt]
      \item \textbf{Debrief item~2 aloud} and leave the four numbers on the board all period;
            rows 1, 2, and~4 all hang on them.
      \item Item~2 is the tell. Any $Q_1$ other than $23$ (usually $22$ or $22.5$) means the
            halves were split wrong --- with ten values each half is five, and $Q_1$ is its
            middle value. Fix it here, once, with the halves on the board.
      \item Item~1's \emph{unsorted} list is deliberate: a student who takes the middle of the
            raw list (``$55$ and $27$'') has not ordered. Say so, briefly.
    \end{itemize}
\end{minipage}
\end{tabularx}
\end{fixedskillbox}

\boxguard
\begin{skillbox}[Hook --- before anyone sits down]{frost}
The hook is on \textbf{page 1 of the notes}, under the vocabulary box, and on the board:
\textit{``Pine Street Pizza's ten Friday deliveries: $20, 22, 23, 24, 25, 26, 27, 28, 30, 55$.
The $55$ was a flat tire. The manager throws it out.''} \textbf{Ask: which number drops more
--- the mean or the median --- or do they drop about the same?} Students circle one and write a
line of reason \emph{before} anyone speaks; then take the hand vote and write the three-way
split on the board. Expect a large ``about the same'' pile --- one value left, so everything
shifts --- and a few who guess the median ``because it's the middle.''
\textbf{Do not settle it.} Say only that the answer is problem~17, in the \emph{Resistance}
row, and that they will vote again there. The vote is what makes the crux land: students have
to have committed to \emph{every summary moves the same} before the computation takes it away.
\end{skillbox}

\boxguard
\begin{skillbox}[Guided Notes \& Practice --- I do, then we do (45 min)]{frost}
\textbf{This is the whole lesson.} There is no group activity, and there is no independent
practice set on this page: \textbf{the homework is the individual practice}, assigned at the
close. The notes are \textbf{one Main Ideas / Notes table}: four instruction rows (problems
1--18), then the Guided Practice row (19--22). \textbf{The rhythm in every row is the same}:
read the definition aloud, mark up the display, work the first problem of the grid yourself,
then hand the rest to the room with the pen in their hand. The vocabulary box on page~1 is
filled \emph{as each term is named}, never in advance. There are no blanks in the notes ---
answers go in each problem's own space and in the \emph{labelboxes} under the displays.
\begin{multicols}{2}
\textbf{I do / we do, row by row (about 25 min).}

\emph{Row 1, Boxplot (problems 1--4), 5 min.} Read the definition --- including the
parenthetical halves rule, the only place it is printed. Point at the four sections of the
Friday boxplot and fill the labelboxes: $25\%$ each, $50\%$ in the box. Work problem~1 yourself
(the warm-up's numbers, read back off the picture). Problem~2 is the ``long whisker holds more''
reflex --- nine times as long, the same quarter of the data. Problem~4 is \emph{finishing marks
on a printed axis}, not a sketch: box $22$--$27$, line at $24$, whiskers to $18$ and~$33$.

\emph{Row 2, Outliers (problems 5--8), 6 min.} Read the four steps; name the dashed lines
(\emph{fences}) and the dot (\emph{outlier}). Work problem~5 live: IQR $5$, step $7.5$, fences
$15.5$ and $35.5$. \textbf{Take problem~7 as a hand vote before anyone writes}: a $16$-minute
delivery that \emph{looks} far. Most say outlier; the fence says $16 > 15.5$. That is the
secondary misconception, and it is settled by arithmetic, not by you. Problem~8 is why the
modified whisker stops at $30$.

\emph{Row 3, Cumulative Graph (problems 9--12), 5 min.} Read the definition; trace
\emph{across from $50\%$, then down} with your finger and fill the labelbox ($25$ minutes).
Work problem~9 ($Q_1 = 20$, $Q_3 = 30$, IQR $10$). Problems 10--12 are percentile language:
$35$ minutes is the $90$th percentile, so $10\%$ of $40 = 4$ free pizzas; Tomas's $30$ minutes
is the $75$th.

\textbf{Row 4, Resistance (problems 13--18) --- this row is the lesson}, 9 min. Put the two dot
plots up. Work problem~13 live (mean $28 \to 25$); the class does 14 (median $25.5 \to 25$) and
15 (IQR $5 \to 5$, with the halves printed for them); 16 hands them the SDs ($9.9 \to 3.1$) to
sort. \textbf{Then re-take the hook vote at problem~17 before you say anything} --- the numbers
are on their page now. Land the printed box word by word: \emph{the median and IQR are
resistant; the mean and standard deviation are not.} Problem~18 is the decision the unit
cares about: the ad's ``average $28$'' describes no typical Friday delivery.

\columnbreak
\textbf{We do (about 20 min).} The Guided Practice row, \emph{Shelter Stays} (problems 19--22)
--- the same moves the homework will ask alone, on a shelter instead of a pizza shop, with an
\emph{even} count so the halves split cleanly. \textbf{Students hold the pen; you hold the
questions.} Problem~19 goes on them cold (five-number summary $2, 4, 6.5, 10, 37$; IQR $6$).
Problem~20: fences $-5$ and $19$ --- a negative lower fence is fine, it just means no low
outliers; the whisker stops at $15$. \textbf{Problem~21 is the crux transferred}: take a hand
count on \emph{mean or median} before anyone computes, then mean $10 \to 7$, median
$6.5 \to 6$. Problem~22 is the report: the flyer's ``$10$ days on average'' should be the
median, $6.5$ days, with the IQR, $6$ days.

Circulate during Guided Practice and demand the reason, not the word:
\begin{itemize}[leftmargin=1.1em, nosep]
  \item \textbf{Problem 19:} ``Show me your two halves. Where did the median go?'' (Even count:
        the halves are the lower five and the upper five.)
  \item \textbf{Problem 20:} ``Is $15$ an outlier? It's far from $10$.'' Make them say the fence.
  \item \textbf{Problem 21:} ``Why did the median barely move when you threw out the biggest
        number?'' A student who says ``it did move the same'' has not left the hook vote.
  \item \textbf{Problem 22:} ``How many of the ten dogs waited $7$ days or fewer? Then is $10$
        a typical wait?'' (Six of ten.)
\end{itemize}
While circulating, pick the papers for the debrief --- one problem~21 that computed both and
said \emph{the mean moved more}, and one that moved both by the same amount. \textbf{This
circulation is your only formative read} --- the homework is not started in class.
\end{multicols}
\end{skillbox}

\boxguard[24]
\begin{skillbox}[Debrief --- whole class, spoken (5 min)]{frost}
Pens down, papers up. \textbf{This debrief is spoken} --- there is no box at the end of the
notes to fill in. \textbf{Five minutes: one board item and the cold check, nothing more.}
\begin{multicols}{2}
\textbf{The board item.} Put the with/without table back up --- mean $28 \to 25$, median
$25.5 \to 25$, IQR $5 \to 5$, SD $9.9 \to 3.1$ --- and have a student say the sentence:
\emph{the median and IQR are resistant; the mean and SD are not}. Then ask the one follow-up
that names the reason: \emph{``What does the mean use that the median doesn't?''} (The value
$55$ itself; the median only uses the middle position.) Hold up the two problem~21 papers you
picked only if there is time --- the cold check comes first.

\columnbreak
\textbf{Check for understanding --- ask it cold, whole class.} \emph{``Ten quiz scores have a
mean of $78$ and a median of $81$. The teacher finds one score was typed as $8$ instead of
$80$ and fixes it. Which changes more --- the mean or the median?''} Hands down, thirty seconds
of think time, then cold-call three students. Correct answer: \textbf{the mean} --- it rises
$72 \div 10 = 7.2$ points, to $85.2$; the median stays about $81$, because fixing the lowest
score does not change the middle. \emph{This replaces the exit ticket.}

\textbf{Formative read.} Sort the answers into three piles: \textbf{(a)} the mean, with the
reason (it uses every value); \textbf{(b)} the mean, no reason, or ``the median doesn't change''
with no why; \textbf{(c)} ``both change the same'' --- the misconception survived. \textbf{If
pile (c) runs past a third of the room, open Lesson~2.4 by re-running problems 13--17}, and say
so plainly.
\end{multicols}
\end{skillbox}

\boxguard
\begin{skillbox}[Homework --- scored, assigned today, due after two study halls]{goldbox}
Six items on \textbf{Elm Park Bike Shop} (eleven used bikes, \$40 to \$900 --- an \emph{odd}
count, so the halves rule is exercised the other way from the notes): \textbf{1} the
five-number summary, with the halves shown (\emph{the core procedure}); \textbf{2} the fences
($30$ and $270$), the outlier (\$900), and whether the \$40 bike that ``looks far'' is one (it
is not --- the \emph{secondary misconception}); \textbf{3} the \textbf{contrast pair}: mean and
median with and without the \$900 bike (\$210 $\to$ \$141 against \$150 $\to$ \$145);
\textbf{4} \emph{the crux head-on}: a classmate's claim that the median drops as much as the
mean, then naming which summaries resist (IQR \$60 $\to$ \$60, $s$ \$233 $\to$ \$46);
\textbf{5} interpret a cumulative relative frequency graph of the season's $40$ bikes (median
\$150, $90\%$ at or below \$250, \$200 is the $75$th percentile); \textbf{6} \emph{the spiral}
to Lesson~2.2 --- the range, \$860 $\to$ \$160, is not resistant, so report the IQR.

\textbf{Assigned at Close \& Assign and done outside class.} \textbf{Scoring:} 12 points ---
2 per item. \textbf{Item~4 is where the misconception shows}; item~3 is where it first surfaces.

\textbf{Packet, Desmos, or DeltaMath --- decided here: DeltaMath + justifications.} The packet
pages are authored and are the default; no Desmos activity asks for the written reasons this
lesson exists for. Items 1, 3, and~5 are the computations and the graph read, which DeltaMath
drills; items 2, 4, and~6 are the ones the teacher note already scores for the \emph{reason}. So
assign \textbf{DeltaMath} --- Box-and-Whisker Plots (five-number summary) $\cdot$ Outliers with
the $1.5 \times \text{IQR}$ rule $\cdot$ Effect of Outliers on Mean and Median $\cdot$ Cumulative
Relative Frequency --- for the computation (items 1, 3, and~5 stay in the packet as a worked
reference), and students \textbf{turn in packet items 2, 4, and~6}. Still 12 points: 2 per
turned-in item, the DeltaMath score scaled to the other 6. Say so aloud at Close \& Assign.

\textbf{Preview of Lesson~2.4.} The spiral box ends on two boxplots on one axis: Elm Park beside
another shop. Do not draw it here --- it is 2.4's opening.

\textbf{Connections \& Big Ideas.} \textit{Unit 2 core idea:} a summary is a claim about the
data, and the right one depends on the data's shape. \textit{Standards covered:} PS.DS.1a
(boxplots and cumulative relative frequency graphs), PS.DS.2c (the $1.5 \times \text{IQR}$
outlier rule), PS.DS.2d (the influence of outliers --- resistance). \textit{Spirals forward
to:} parallel boxplots and comparing distributions (2.4, PS.DS.1a), and the choice of center
and spread every data project in the course must defend. \textit{Reaches back to:} Lesson~2.2's
mean, median, quartiles, IQR, range, and standard deviation, and Algebra~1 percent reasoning.
\end{skillbox}

\boxguard
\begin{skillbox}[Watch For (while circulating)]{redbox}
\begin{itemize}[leftmargin=1.2em, nosep]
  \item \textbf{``The outlier moves every summary the same''} (problems 17--18, GP~21--22,
        HW~3--4 --- the target misconception). Probe: \emph{``Which of the ten values does the
        median actually use? Did the $55$ change which value is in the middle?''}
  \item \textbf{``It looks far, so it's an outlier''} (problem~7, GP~20, HW~2 --- the second
        misconception). Probe: \emph{``What is the fence? Which side of it is the point on?''}
  \item \textbf{The modified whisker drawn to the outlier} (problem~8, GP~20, HW~2). Probe:
        \emph{``The whisker reaches the biggest value that is \emph{not} an outlier. Which one
        is that?''}
  \item \textbf{A long whisker holds more data} (problem~2). Probe: \emph{``How many deliveries
        are in the right whisker? In the left?''}
  \item \textbf{The halves done wrong} (warm-up~2, GP~19, HW~1): the median put into a half when
        $n$ is odd, or the list never sorted. Probe: \emph{``Show me your lower half. How many
        values are in it?''}
  \item \textbf{The cumulative graph read down-then-across, or the percent read as a time}
        (problems 9--12, HW~5). Probe: \emph{``Start on the percent axis. Across to the curve ---
        then where?''}
  \item \textbf{The SD called resistant ``because it's spread, like the IQR''} (problem~16,
        HW~4). Probe: \emph{``Does the SD use the $55$ in its arithmetic?''}
\end{itemize}
Cold-call prompts: \emph{``Give me a data set from this school where the mean would mislead.''}
\quad \emph{``Name a summary from Lesson~2.2 that is \emph{not} resistant besides the mean.''}
(The range, and the SD.)
\end{skillbox}

\boxguard
\begin{skillbox}[Close \& Assign (5 min)]{goldbox}
\textbf{Launch the homework --- it is assigned now and done outside class, not started here.}
Say the form (\textbf{DeltaMath} for the computation, and \textbf{packet items 2, 4,
and~6 turned in}, pages after the notes), the points (\textbf{12}, 2 per
item), and the due date (\textbf{the start of the first class after two study halls}); point
out that the \emph{Keep in Mind} box on the cover is the page to flip back to. Tell them to do
\textbf{item~1 first} --- every later item leans on its quartiles --- and that \textbf{item~4 is
the one you will read first}: a student who says the median drops as much as the mean has
learned the words and not the rule. Then close on what changed: \emph{``You walked in able to
compute a median and a mean. You walk out knowing that one flat tire moves the mean three
minutes and the median half a minute --- and that the fence, not your eye, decides what an
outlier is.''} \textbf{Preview:} Lesson~2.4, \emph{Comparing Distributions} --- two boxplots
on one axis, and a claim about which group is ``faster'' or ``cheaper'' that you will have to
defend with the resistant summaries.
\end{skillbox}

% ── Teacher notes — one per component, in packet order (three) ──────────────
% TEACHERNOTES: teacher prose lives HERE, never in a _key. No note for the debrief or the
% homework launch — both are fully specified in their own boxes above.
\begin{teachernote}[Warm-Up]
Five minutes, silent, timed. This is Lesson~2.2's arithmetic on the exact numbers the notes open
with --- say so, so nobody thinks they are being retested. \textbf{Item~2 is the one to read
while collecting}: the halves method is the single most common error of the unit, and
everything in rows 1, 2, and~4 is built on $Q_1 = 23$ and $Q_3 = 28$. Any other $Q_1$ means
the halves were split wrong: ten values make two halves of five, and each quartile is the
middle of its five. Put the ordered list and the two halves on the board, write
$Q_1 = 23$, median $25.5$, $Q_3 = 28$, IQR $5$ beside them, and \textbf{leave them up all
period}. Item~3's three percents are the three dots of row~3's graph at $25\%$, $50\%$, $75\%$;
point that out when you get there, not now.
\end{teachernote}

\begin{teachernote}[Guided Notes \& Practice]
\textbf{This one document is the lesson --- pace it 25 / 20, then 5 for the debrief and 5 to
assign.} It is one table, and every row runs the same way: you read the definition, mark up the
display, and work the first problem of the grid (1, 5, 9, 13); the class works the rest with
the pen in their hand. Rows 1--3 must move --- five, six, and five minutes. \textbf{Take the
hand count on problem~7 before anyone writes}; if they compute first, the ``looks far'' trap is
gone. Problem~4 finishes a boxplot on a printed axis --- walk the room once to check that the
box spans $22$ to $27$, not $18$ to $33$.

\textbf{Spend the time in row~4, \emph{Resistance}}, which carries the misconception --- nine
minutes. Work problem~13 live, give 14 and~15 to the room (the halves for 15 are printed so it
stays quick), hand them the SDs in 16 to sort, and \textbf{re-take the hook vote at problem~17
before you say anything}. Let the room see its own numbers disagree with its vote: the mean
moved six times as far as the median. Only then read the printed box. Do not resolve it by
restating the definition; put \emph{``$3$ minutes''} and \emph{``$0.5$ minute''} side by side
on the board and ask what the mean uses that the median does not.

\textbf{Guided Practice (19--22) is where the pen changes hands.} \emph{Shelter Stays} runs
the four moves the homework asks alone --- summary, fences, with and without, the report ---
on a shelter instead of a pizza shop. Problem~21 is the crux transferred; a student who moves
both summaries by the same amount has not left the hook vote. \textbf{If time is short,
protect 21 and~22's verdict}; problem~20's whisker question can be said aloud.

\textbf{There is no independent practice set on this page, and the homework is not started in
class} (the Unit~2 period). Your formative read is the Guided Practice circulation and the
debrief's cold check. If the ``both move the same'' pile is more than a third of the room,
\textbf{open Lesson~2.4 by re-running problems 13--17}, and say so plainly.
\end{teachernote}

\begin{teachernote}[Homework]
\textbf{Scored, 12 points (2 per item), due at the start of the first class after two study
halls; assigned at Close \& Assign, done outside class.} Score items 1, 3, and~5 for
correctness (one point per half; accept \$145 or ``halfway between \$140 and \$150'' for the
median without the \$900 bike). Score items 2, 4, and~6 for the \emph{reason}: item~2 needs a
fence named for \emph{both} the \$900 and the \$40 bike; item~4 needs the two changes compared
(\$69 against \$5) \emph{and} the four summaries sorted; item~6 needs ``it uses the \$900
itself'' or equivalent, not just ``not resistant.'' \textbf{Read item~4 first, and sort it into
three piles} --- the median barely moved, with a reason (ready); the right verdict, no reason
(ready, needs the language); ``she is right'' or both summaries moved equally (the misconception
survived). If more than a third land in the last pile, reopen the Pine Street with/without
table in the Lesson~2.4 warm-up debrief rather than letting it ride.

\textbf{DeltaMath + justifications is the call for 2.3.} DeltaMath: Box-and-Whisker Plots
(five-number summary) $\cdot$ Outliers with the $1.5 \times \text{IQR}$ rule $\cdot$ Effect of
Outliers on Mean and Median $\cdot$ Cumulative Relative Frequency. Students turn in packet items
2, 4, and~6. Score items 2, 4, and~6 as above (item~4 first); scale the DeltaMath score to the
other 6 points. Students still work item~1 on the page as scratch --- item~2 needs its quartiles.
Confirm the topic names against DeltaMath's Statistics module when you build the assignment.
Item~5's graph read is now on DeltaMath --- if the cumulative-frequency problems come back weak,
start 2.4's warm-up with one ``across, then down'' read.
\end{teachernote}
\end{document}
